\chapter{Ứng dụng Flutter}

\section{Tổng quan}

\begin{longtable}{L{4.4cm} L{9.2cm}}
\toprule
\textbf{Thuộc tính} & \textbf{Giá trị}\\
\midrule
\endhead
Phiên bản & 1.7.0+17\\
Dart SDK & \(\ge\)3.8.1\\
Nền tảng & Android, iOS, Web (kèm macOS, Linux, Windows ở mức thử nghiệm)\\
Quản lý trạng thái & Provider (\code{ChangeNotifier}) — 31 provider đăng ký ở \code{main.dart}\\
Tiêm phụ thuộc & GetIt, khởi tạo theo mô-đun tính năng\\
Kiến trúc & Clean Architecture cho từng tính năng\\
Lưu trữ cục bộ & \code{sqflite} (di động), \code{sqflite\_ffi\_web} (web), \code{flutter\_secure\_storage} (thẻ)\\
Đa ngữ & \code{easy\_localization}, tệp JSON trong \code{assets/i18n/}\\
Số mô-đun tính năng & 26\\
Số màn hình & khoảng 90\\
Số gói phụ thuộc & 63\\
Số tệp kiểm thử & 120\\
\bottomrule
\caption{Thông số ứng dụng Flutter}
\end{longtable}

\section{Kiến trúc sạch theo tính năng}

Mỗi tính năng trong \code{lib/features/\{ten\}/} có cùng một hình dạng:

\begin{lstlisting}
{feature}/
  data/
    datasources/     # nguon xa (HTTP/WS) va nguon cuc bo (SQLite)
    models/          # model JSON, fromJson/toJson
    repositories/    # cai dat cua interface o domain
  domain/
    entities/        # thuc the Dart thuan, khong phu thuoc gi
    repositories/    # interface truu tuong
    usecases/        # moi use case mot trach nhiem
  presentation/
    screens/ | pages/
    providers/       # ChangeNotifier
    widgets/
  di/
    {feature}_module.dart   # dang ky GetIt
\end{lstlisting}

Luật bất di bất dịch: \code{domain/} \textbf{không được} nhập bất cứ thứ gì từ
\code{presentation/} hay \code{data/}. Chiều phụ thuộc luôn hướng vào trong.

Luồng dữ liệu một lần chạm màn hình:

\begin{center}
\code{Widget $\rightarrow$ Provider $\rightarrow$ UseCase $\rightarrow$ Repository $\rightarrow$
RemoteDataSource $\rightarrow$ Gateway}
\end{center}

với \code{LocalDataSource} làm bộ đệm song song ở tầng repository.

\section{Hai mươi sáu mô-đun tính năng}

\begin{longtable}{L{3.2cm} L{10.4cm}}
\toprule
\textbf{Mô-đun} & \textbf{Nội dung}\\
\midrule
\endhead
\code{auth} & Đăng nhập email, Google, Facebook; đăng ký; đặt lại mật khẩu; lớp bảo vệ tuyến.\\
\code{home} & Bảng điều khiển chính, tổng hợp dữ liệu nhiều nguồn.\\
\code{learning} & Lộ trình học kiểu Duolingo: các nút zigzag vẽ trên \code{CustomPaint} với đường Bézier bậc ba, nút bài học hiện tại có hoạt ảnh nhấp nháy.\\
\code{course} & Danh mục khoá học theo bậc CEFR.\\
\code{vocabulary} & Thẻ từ, ôn tập theo lịch FSRS, bộ thẻ tự tạo.\\
\code{practice} & Phòng luyện tập tổng hợp.\\
\code{mistakes} & Sổ tay lỗi sai, đánh dấu đã ôn hoặc mở lại.\\
\code{lexi\_chat} & Trợ lý Lexi chạy TRACE-CAG, nhận phản hồi theo dòng.\\
\code{chat} & Hội thoại theo chủ đề.\\
\code{voice} & Luyện nói thời gian thực, phản hồi âm vị.\\
\code{level} & Theo dõi bậc CEFR và hồ sơ năng lực.\\
\code{progress} & Thống kê, chuỗi ngày, bản đồ nhiệt hoạt động.\\
\code{games} & Sáu trò chơi mini.\\
\code{gamification} & XP, ví, cửa hàng, kho đồ, bảng xếp hạng.\\
\code{achievements} & Hệ thống huy hiệu.\\
\code{social} & Theo dõi bạn bè, bảng tin hoạt động.\\
\code{books}, \code{news}, \code{podcast}, \code{youtube} & Bốn loại nội dung thực tế kèm quiz.\\
\code{notifications} & Thông báo trong ứng dụng và đẩy qua FCM.\\
\code{premium} & Gói trả phí, quyền lợi.\\
\code{offline} & Chế độ ngoại tuyến và đồng bộ lại.\\
\code{profile}, \code{user} & Hồ sơ, cài đặt, quản lý tài khoản.\\
\code{admin} & Màn hình quản trị nhúng cho vai trò quản trị viên.\\
\bottomrule
\caption{Mô-đun tính năng của ứng dụng Flutter}
\end{longtable}

\section{Trình tự khởi động}

\begin{lstlisting}[language=Java]
main() async
  EasyLocalization.ensureInitialized()
  LocaleService.getSavedLocale()
  databaseFactory = ffiWebNoWebWorker        // chi tren web
  dotenv.load(.env | .env.production)
  Firebase.initializeApp()
  LocalStateMigrationService().runIfNeeded()  // di tru du lieu cuc bo
  di.initializeDependencies(skipDatabase: kIsWeb)
  NotificationService.ensureInitialized()
  StartupCoordinator.run()                    // kiem tra suc khoe backend
  runApp(EasyLocalization -> LexiLingoApp)

// sau khung hinh dau tien
WidgetsBinding.addPostFrameCallback:
  FirebaseMessagingService.initialize()       // xin quyen + dang ky token
\end{lstlisting}

Việc khởi tạo FCM được hoãn tới sau khung hình đầu tiên vì hộp thoại xin quyền thông báo
sẽ chặn giao diện nếu bật ngay lúc khởi động.

\section{Dịch vụ lõi}

\begin{longtable}{L{4.8cm} L{8.8cm}}
\toprule
\textbf{Dịch vụ} & \textbf{Chức năng}\\
\midrule
\endhead
\code{ApiClient} & Máy khách HTTP có ba bộ chặn: gắn thẻ xác thực, thử lại khi 401 kèm làm mới thẻ, ghi nhật ký gỡ lỗi.\\
\code{skill\_event\_recorder} & Gửi sự kiện học tập lên backend theo kiểu bắn-và-quên.\\
\code{sync\_queue\_lifecycle\_runner}, \code{background\_sync\_queue\_service} & Hàng đợi hành động ngoại tuyến; đồng bộ khi ứng dụng quay lại tiền cảnh.\\
\code{progress\_sync\_service}, \code{progress\_firestore\_data\_source} & Đồng bộ tiến độ, có đường dự phòng Firestore.\\
\code{encrypted\_local\_cache\_service}, \code{local\_cache\_service} & Bộ đệm cục bộ, có mã hoá cho dữ liệu nhạy cảm.\\
\code{user\_scope\_service} & Cô lập dữ liệu cục bộ theo từng tài khoản, tránh rò dữ liệu giữa hai lần đăng nhập khác nhau.\\
\code{session\_expired\_service} & Xử lý tập trung khi phiên hết hạn.\\
\code{deep\_link\_service} & Liên kết sâu, ví dụ đặt lại mật khẩu qua email.\\
\code{podcast\_audio\_handler} & Phát âm thanh nền qua \code{audio\_service}.\\
\code{known\_words\_service}, \code{quick\_save\_vocabulary\_service} & Đánh dấu từ đã biết, lưu nhanh từ mới khi đang đọc.\\
\code{purchases\_service}, \code{entitlement\_service} & Mua trong ứng dụng và quyền lợi gói.\\
\code{firebase\_messaging\_service}, \code{notification\_service} & Thông báo đẩy và thông báo cục bộ.\\
\code{health\_check\_service} & Kiểm tra backend còn sống trước khi vào ứng dụng.\\
\code{sound\_service}, \code{rating\_service}, \code{analytics\_service} & Âm thanh phản hồi, nhắc đánh giá ứng dụng, đo lường.\\
\bottomrule
\caption{Dịch vụ lõi phía client}
\end{longtable}

\section{Hoạt động ngoại tuyến}

Ứng dụng phải dùng được khi mạng chập chờn. Cơ chế gồm ba phần:

\begin{enumerate}[leftmargin=1.4em]
  \item \textbf{Bộ đệm đọc}: dữ liệu khoá học, từ vựng và tiến độ được lưu vào SQLite;
  màn hình đọc từ bộ đệm trước rồi làm mới nền.
  \item \textbf{Hàng đợi ghi}: mọi hành động thay đổi dữ liệu khi ngoại tuyến được xếp vào
  \code{sync\_queue} cục bộ.
  \item \textbf{Đồng bộ theo vòng đời}: \code{SyncQueueLifecycleRunner} chạy khi ứng dụng
  quay lại tiền cảnh (\code{didChangeAppLifecycleState}) và theo chu kỳ, đẩy hàng đợi lên
  máy chủ.
\end{enumerate}

Trên nền web, SQLite dùng bản \code{ffi} không có \en{web worker}, nên một số bước bị bỏ
qua có chủ đích: nạp hồ sơ XP và danh sách podcast tuyển chọn được hoãn lại, kiểm tra sức
khoẻ backend bị bỏ.

\section{Giao diện và trải nghiệm}

\begin{itemize}[leftmargin=1.4em]
  \item \textbf{Chủ đề}: \code{AppTheme.lightTheme} / \code{darkTheme}; mọi widget lấy màu
  và cỡ chữ từ \code{Theme.of(context)}, không mã hoá cứng giá trị.
  \item \textbf{Lộ trình học}: các nút được đặt bằng \code{Positioned} trong một
  \code{Stack}, toạ độ ngang theo chu kỳ hình sin
  \([0{,}14;\ 0{,}32;\ 0{,}50;\ 0{,}68;\ 0{,}86;\ 0{,}68;\ 0{,}50;\ 0{,}32]\);
  \code{RoadmapPathPainter} vẽ đường Bézier bậc ba với ba lớp (bóng, đường chính, nét đứt
  làm điểm nhấn).
  \item \textbf{Cuộn}: \code{BouncingScrollPhysics} trên iOS, \code{ClampingScrollPhysics}
  trên Android để hợp quy ước từng nền tảng.
  \item \textbf{Trạng thái tải}: \code{shimmer} làm khung xương, \code{lottie} cho hoạt ảnh
  thành tích, \code{confetti} khi hoàn thành mục tiêu.
\end{itemize}

\section{Đa ngôn ngữ}

Tệp dịch nằm ở \code{assets/i18n/*.json}, ngôn ngữ dự phòng là tiếng Anh. Khoá dịch được
đồng bộ với Crowdin qua quy trình CI (\code{crowdin-sync.yml}), và máy chủ MCP cung cấp
công cụ \code{manage\_i18n\_key} để thêm hoặc sửa một khoá trên tất cả tệp ngôn ngữ cùng
lúc — tránh tình trạng khoá tồn tại ở tệp này nhưng thiếu ở tệp khác.
